% !TEX program = pdflatex
% ------------------------------------------------------------
% MATH 521 Homework Template (adapted from your MATH 421 HW10)
%
% HOW TO USE (quick):
%   1) Edit the Metadata block (HW number / due date especially).
%   2) Paste each problem statement under \problem[]{<n>}.
%   3) Write in the Scratch/Discussion box if you want.
%   4) Write your proof under "Proof." and end with \qed (or \qedhere).
%
% Notes:
% - Keeps theorem/definition boxes (tcolorbox).
% - Keeps ToC + PDF bookmarks for each problem.
% - Adds a Scratch/Discussion box per problem (optional).
% - Header bar has course + meeting info (left) and HW number (right).
% ------------------------------------------------------------

% TROUBLESHOOTING NOTE:
% If you ever see errors like `\tcb@insert@after@part` or other tcolorbox internals being
% an "Undefined control sequence", that usually means LaTeX is accidentally loading a
% DIFFERENT (often older) `tcolorbox.sty` from your project folder or local texmf tree.
% Fix: search your repo for `tcolorbox.sty` and delete/rename any local copy so TeX uses
% the TeX Live version. Also try "Clean" in LaTeX Workshop to remove stale .aux/.toc/.out.

\documentclass[12pt]{article}

% =========================
% 0) Metadata (EDIT ME)
% =========================
\newcommand{\CourseCode}{MATH 521}
\newcommand{\CourseTitle}{Analysis I}
\newcommand{\CourseSection}{LEC 002}
\newcommand{\Semester}{Spring 2026}
\newcommand{\Instructor}{Thierry Laurens}
\newcommand{\MeetingInfo}{MWF 1:20--2:10 \;|\; Van Vleck B119} % optional
\newcommand{\HomeworkNumber}{3} % <-- change each week
\newcommand{\YourName}{Alejandro De La Torre}
\newcommand{\DueDate}{February 13, 2026} % <-- or set e.g. January 31, 2026

% =========================
% 1) Core math + proof tools
% =========================
\usepackage{amsmath, amssymb, amsthm}

% =========================
% 2) Lists and spacing
% =========================
\usepackage{enumitem}
\setlist[enumerate,1]{label=\textbf{(\alph*)}, leftmargin=2em, itemsep=0.25em}
\setlist[itemize]{leftmargin=1.5em, itemsep=0.25em}
\setlength{\parskip}{0.35em}
\setlength{\parindent}{0pt}

% =========================
% 3) Page geometry + header/footer
% =========================
\usepackage{geometry}
\geometry{margin=1in}

\usepackage{fancyhdr}
\pagestyle{fancy}
\fancyhf{}
\lhead{\CourseCode\ --- \CourseSection, \Semester \;(\MeetingInfo)}
\rhead{Homework \HomeworkNumber}
\cfoot{\thepage}
\renewcommand{\headrulewidth}{0.4pt}
% Fixes common fancyhdr warning about header height:
\setlength{\headheight}{15pt}
% Optional: reclaim a bit of vertical space after increasing headheight
\addtolength{\topmargin}{-2pt}

% =========================
% 4) Links (subtle)
% =========================
\usepackage[hidelinks]{hyperref}
\setcounter{tocdepth}{1}

% =========================
% 5) Quotes / figures
% =========================
\usepackage{csquotes}
\usepackage{graphicx}
\graphicspath{{images/}} % put figures in ./images/

% =========================
% 6) Simple scratch box (no tcolorbox)
% =========================
% A lightweight environment for brief planning / scratch work.
% This avoids any tcolorbox dependencies.
\newenvironment{scratch}{%
  \par\smallskip
  \noindent\textbf{Discussion \& Scratch Work}\begin{quote}\small
}{%
  \end{quote}\smallskip
}

% =========================
% 7) Lightweight theorem-like environments (optional)
% =========================
\newtheorem*{claim}{Claim}
\newtheorem*{lemma}{Lemma}
\newtheorem*{remark}{Remark}
\newtheorem{theorem}{Theorem}
\newtheorem{proposition}{Proposition}
\newtheorem{corollary}{Corollary}
\newtheorem{definition}{Definition}
\newtheorem{example}{Example}
\newtheorem{exercise}{Exercise}

% =========================
% 8) Problem header command (with TOC + PDF bookmarks)
% Usage: \problem[Optional short title]{<number>}
% =========================
\makeatletter
\newcommand{\problem}[2][]{%
  \def\prob@title{Problem #2\if\relax\detokenize{#1}\relax\else\ (\textit{#1})\fi}%
  \phantomsection%
  \pdfbookmark[1]{\prob@title}{prob#2}%
  \addcontentsline{toc}{section}{\prob@title}%
  \section*{\prob@title}%
  \vspace{-0.25em}\hrule\vspace{0.8em}%
}
\makeatother

% =========================
% 9) Title block
% =========================
\title{Homework \HomeworkNumber}
\author{\YourName \\
\CourseCode\ --- \CourseTitle \\
\CourseSection \\
Instructor: \Instructor}
\date{\DueDate}

\begin{document}
\maketitle

% ------------------ collaboration + sources ------------------
% \section*{Collaboration Statement}
% Write “I worked alone” or list collaborators / external resources.
% "I worked on this homework independently. No outside collaborators or resources were used."

\tableofcontents
\bigskip
\hrule
\bigskip

\newpage

% ============================================================
% Problems (duplicate blocks as needed)
% ============================================================

\problem[]{1}
Let $(F, +, \cdot)$ be a field.
\begin{enumerate}
\item Prove that the multiplicative identity is unique: If $b \in F$ satisfies $b \cdot a = a \cdot b = a$ for all $a \in F$, then $b = 1$.
\item Prove that multiplicative inverses are unique: If $a \in F \setminus \{0\}$ and $b \in F$ satisfy $b \cdot a = a \cdot b = 1$, then $b = a^{-1}$.
\item Prove that if $a \cdot b = a \cdot c$ and $a \neq 0$, then $b = c$.
\item Prove that if $a, b \in F$ then $-a + (-b) = -(a + b)$.
\item Prove that if $a, b \in F$ and $a, b \neq 0$, then $a^{-1} \cdot b^{-1} = (ab)^{-1}$.
\end{enumerate}

\begin{scratch}
% Outline the strategy, key definitions, lemmas, or a proof sketch.
\end{scratch}

\textbf{Proof.}\\
% Write your proof here.
\qed

\newpage

\problem[]{2}
Given two points $(x_1, y_1)$ and $(x_2, y_2)$ in $\mathbb{R}^2 = \{(x,y) : x,y \in \mathbb{R}\}$, we write
\[
(x_1, y_1) < (x_2, y_2) \Leftrightarrow x_1 < x_2 \text{ or } (x_1 = x_2 \text{ and } y_1 < y_2).
\]
\begin{enumerate}
\item Prove that $<$ is an order relation on $\mathbb{R}^2$.
\item Prove that $(\mathbb{R}^2, <)$ does not satisfy the least upper bound property. (Hint: Consider the set of all points on the y-axis.)
\end{enumerate}

\begin{scratch}
\end{scratch}

\textbf{Proof.}\\
\qed

\newpage

\problem[]{3}
Recall that the set of complex numbers $\mathbb{C}$ consists of all numbers of the form $a + ib$ for $a, b \in \mathbb{R}$. We can add and multiply complex numbers as follows:
\[
(a + ib) + (c + id) = (a + c) + i(b + d) \quad \text{and} \quad (a + ib)\cdot(c + id) = (ac - bd) + i(ad + bc).
\]
With these operations, it turns out $(\mathbb{C}, +, \cdot)$ is a field (but you do not need to prove this).
\begin{enumerate}
\item The definition of a field says that for any $a + ib \neq 0 + i0$, there exists $c + id \in \mathbb{C}$ so that $(a + ib)\cdot(c + id) = 1 + i0$. What are $c, d \in \mathbb{R}$ in terms of $a, b \in \mathbb{R}$? Prove your answer.
\item Prove that there is no order relation $<$ on $\mathbb{C}$ that makes $(\mathbb{C}, +, \cdot, <)$ an ordered field. (Hint: Can $i = 0 + i1$ be positive? Negative? Zero?)
\end{enumerate}

\begin{scratch}
\end{scratch}

\textbf{Proof.}\\
\qed

\newpage

\problem[]{4}
Let $A \subseteq \mathbb{R}$ be a nonempty and bounded set, let $c \in \mathbb{R}$ be a number, and define
\[
B = \{c \cdot a : a \in A\}.
\]
\begin{enumerate}
\item Prove that if $c > 0$, then $\sup B = c \cdot \sup A$.
\item Prove that if $c < 0$, then $\sup B = c \cdot \inf A$.
\end{enumerate}

\begin{scratch}
\end{scratch}

\textbf{Proof.}\\
\qed

\newpage

\problem[]{5}
Let $A, B \subseteq \mathbb{R}$ be nonempty and bounded sets, and define
\[
S = \{a + b : a \in A \text{ and } b \in B\}.
\]
Prove that $\sup S = \sup A + \sup B$.

\begin{scratch}
\end{scratch}

\textbf{Proof.}\\
\qed

\newpage

\problem[]{6}
Let $A, B \subseteq \mathbb{R}$ be nonempty and bounded sets. Prove or disprove the following statements:
\begin{enumerate}
\item If $a \le b$ for all $a \in A$ and for all $b \in B$, then $\sup A \le \inf B$.
\item If $a < b$ for all $a \in A$ and for all $b \in B$, then $\sup A < \inf B$.
\end{enumerate}

\begin{scratch}
\end{scratch}

\textbf{Proof.}\\
\qed

% ============================================================
% Optional appendix: keep a personal toolbox of definitions/theorems
% ============================================================
\newpage
\appendix
\section*{Appendix: Toolbox (optional)}
\addcontentsline{toc}{section}{Appendix: Toolbox (optional)}

% Example (using amsthm environments instead of boxed tcolorboxes):
% \begin{definition}
% A metric space is ...
% \end{definition}
%
% \begin{theorem}[Heine--Borel]
% ...
% \end{theorem}

\end{document}
